% 2009 年湖北卷文科第 15 题：样本容量 200 的频率分布直方图。
% 来源 PDF：普通高考/2009/2009湖北文.pdf，SHA-256 f60ba7a5a2fef3449fba9474fa973c616aa52a15926b3a97f0dda6b909f4c861。
% 原页：第 1 页第 2 栏下部，书页码 525。
% 重绘依据：横轴 0—22，0 与 2 之间为省略断轴；纵轴为「频率/组距」，组距 4。
%   五组矩形高：[2,6)=0.02，[6,10)=0.08，[10,14)=0.09，[14,18)=0.03，[18,22)=0.03；
%   0.02+0.08+0.09+0.03+0.03=0.25，乘组距 4 恰为 1，已核对。
%   虚线标出 0.02、0.03、0.08、0.09 四条水平参考线，位置照原页。
%   纵轴标题置于轴顶右侧、0.09 虚线之上，避开 0.09/0.08 刻度字与虚线。
\begin{tikzpicture}[x=0.34cm, y=30cm, semithick,
  every node/.style={inner sep=1.2pt, font=\small}]

  % 水平虚线参考线（画在矩形之下，白色填充的矩形覆盖其内部）
  \draw[dashed] (0,0.09) -- (10,0.09);
  \draw[dashed] (0,0.08) -- (10,0.08);
  \draw[dashed] (10,0.03) -- (22,0.03);
  \draw[dashed] (0,0.02) -- (2,0.02);

  % 各组矩形
  \draw[fill=white] (2,0) rectangle (6,0.02);
  \draw[fill=white] (6,0) rectangle (10,0.08);
  \draw[fill=white] (10,0) rectangle (14,0.09);
  \draw[fill=white] (14,0) rectangle (18,0.03);
  \draw[fill=white] (18,0) rectangle (22,0.03);

  % 坐标轴
  \draw[-{Stealth}] (-0.5,0) -- (23.4,0) node[right, inner sep=3pt] {样本数据};
  \draw[-{Stealth}] (0,-0.002) -- (0,0.1125);
  \node[anchor=west] at (0.45,0.1045) {频率/组距};

  % 断轴符号（0 与 2 之间）
  \draw (0.45,0) -- (0.58,0.0022) -- (0.71,0) -- (0.84,0.0022) -- (0.97,0);

  % 刻度与标签
  \foreach \x in {2,6,10,14,18,22} {
    \draw (\x,0) -- (\x,-0.0012);
    \node[below] at (\x,-0.0018) {$\x$};
  }
  \node[below left=1pt] at (0,0) {$O$};

  % 纵轴刻度标签
  \foreach \h/\t in {0.02/0.02, 0.03/0.03, 0.08/0.08, 0.09/0.09} {
    \node[left, font=\scriptsize] at (0,\h) {\t};
  }
\end{tikzpicture}
